% Generated by writeup/fill_tables.py -- do not hand-edit.
% Needs only booktabs + multirow. No custom macros, no xcolor.
% An em-dash cell = not scheduled; 'running' = sweep has not filled it yet.
\begin{table}[h]
  \centering
  \caption{Planning success rate (\%) over 100 seeds per cell, at a matched
  150-iteration budget. Goals are drawn with \texttt{random\_state} on Reacher
  and Point maze and from the dataset on Push-T, whose random start/goal pairs
  are frequently unconnectable. Superscripts give $n$ where an arm is
  incomplete, and a star marks a figure that is still moving: the arm is
  either mid-sweep or a provisional probe value that the full run will
  replace. An unmarked number is final over 100 seeds.
  \textbf{The FK budget is not uniform across environments}:
  Reacher runs 8 particles $\times$ 8 smoothing samples, while Point maze and
  Push-T run 3 particles with smoothing disabled --- the DINO-WM encoder emits
  196 ViT patch tokens per frame and the rollout is retained for autograd, so
  8 particles exceeds 16\,GB of VRAM there. Those two rows are therefore
  reduced-budget for our method and favour the baselines.}
  \label{tab:planners}
  \begin{tabular}{llcccc}
    \toprule
    Environment & $H$ & GD & GRASP & CEM & FK (ours) \\
    \midrule
    \multirow{3}{*}{Reacher}
     & 5  & 57 & 27 & 62 & 100 \\
     & 10 & 38 & 33 & 38 & 92 \\
     & 15 & 16 & 32 & 28 & 50 \\
    \midrule
    \multirow{3}{*}{Point maze}
     & 5  & 54 & 44 & 51 & 75 \\
     & 10 & 51 & 49 & 67 & 81 \\
     & 15 & 50 & 62 & 69 & 61 \\
    \midrule
    \multirow{3}{*}{Push-T}
     & 5  & 57\textsuperscript{81} & \textendash & \textendash & 68 \\
     & 10 & \textendash & \textendash & \textendash & 27 \\
     & 15 & \textendash & \textendash & \textendash & 5 \\
    \midrule
    \multirow{3}{*}{Wall}
     & 5  & 28 & 21 & 33 & \textendash \\
     & 10 & 41 & 9 & 40 & \textendash \\
     & 15 & 25 & 5 & 49 & \textendash \\
    \midrule
    \multirow{3}{*}{Cube}
     & 5  & \textendash & \textendash & \textendash & \textendash \\
     & 10 & \textendash & \textendash & \textendash & \textendash \\
     & 15 & \textendash & \textendash & \textendash & \textendash \\
    \bottomrule
  \end{tabular}
\end{table}
